\section{Provenance and Numerical Audit}
\label{sec:credibility_audit}

The preceding experiments answer whether BC-assisted direct TD-MPC2 can complete the
evaluated task and where its behaviour changes. This section asks a different
question: which parts of that evidence retain the same interpretation after
the geometric sources and numerical settings are inspected?

For the 288 matched-OOD evaluations, a separate audit program re-read the
frozen JSONL traces, verified their hashes, and recomputed all 7,776 reported
per-evaluation metric values from an independently implemented specification.
The program imports neither the experimental analysis nor the finalizer or
runner. Every value matched the stored summary, with no trace-hash or
trace-identity mismatch. Here, independent refers to the metric-computation
logic, not to an independently generated trajectory set.

\subsection{Geometric and Disturbance Provenance}

Table~\ref{tab:provenance} separates information available before action
selection from information that would be available through a physical sensing
stack. The controller is time-causal, but it uses simulator state and a
registered nominal surface model. The latter supplies the local frame for both
observations and action execution; it is not inferred by TD-MPC2.

\begin{table*}[t]
\centering
\caption{Sources of quantities used by the controller and evaluator.}
\label{tab:provenance}
\footnotesize
\setlength{\tabcolsep}{3pt}
\begin{tabular}{@{}p{0.18\textwidth}p{0.25\textwidth}p{0.47\textwidth}@{}}
\toprule
Quantity & Source & Role and interpretation \\
\midrule
Normal force and rate & Simulator contact readback & Observation, routing, and force metrics \\
Tool pose and velocity & Simulator state & Surface-relative observation \\
Path tangent and surface normal & Registered nominal model & Observation frame and Cartesian action map \\
Progress and contact dose & Synthetic task state & Observation and completion proxy \\
\bottomrule
\end{tabular}
\end{table*}

A separate claim--binding audit covered all controlled and OOD block labels.
The physical call chains for the designated lateral and normal impulses were
present. By contrast, the audit did not establish that the declared reference
noise and delayed sensor noise in two controlled and two OOD blocks reached
the learner-visible force path. Those labels are therefore not used as
evidence of noise robustness.

\subsection{Outcome-Specific Numerical Sensitivity}

All four 12~N OOD baseline sentinels reproduced their archived traces within the
prespecified tolerance. Supplementary Tables S2 and S3 expose all 32 paired
intervention cells.
Each force entry is the candidate peak in
newtons, followed by three pass bits for task, tracking, and the force limit.
The sentinel headers give the corresponding reproduced-baseline values. An
arrow marks a trigger of the prespecified material peak/threshold rule and its
direction; an unmarked value did not trigger the rule. Status changes relative
to baseline are summarized below from the same paired cells.

The rule triggered in 11/32 cells. Task-success status changed in three cells,
tracking status in two, and sampled force-limit status in two. Ten of the
eleven material peak responses were reductions. The remaining response was an
increase that crossed the 15~N threshold.

Aggregate pass--fail counts alone mask these cell-level outcome switches.
Repeating each block's
baseline outcome across the eight interventions gives 16 task, 16 tracking,
and 16 force-limit passes. The candidates give 15, 16, and 16, respectively,
despite the paired cell-level changes in
the complete cell tables in the Supplementary Material.
The relevant finding is therefore outcome-specific sensitivity, not a binary
division between invariant task results and unstable force results.

All twelve requested source arc counts reached the mesh factory, but PhysX
convex cooking returned baseline-equivalent collision payloads in eleven
cells (Supplementary Material). Only the 96-segment F4 request
changed the topology, from 14 vertices and 24 triangles to 16 and 28, and produced the
arm's sole material peak response. The arm thus records request-to-collision
mapping rather than multilevel mesh convergence. The 0.080~mm analytical
source-mesh sagitta bound makes penetration-depth error unlikely as the sole explanation.
Position iterations and contact offset also affect selected sentinels; physics
substeps remain fixed in this 100~Hz, one-step panel.

\subsection{Resulting Evidence Boundaries}

Table~\ref{tab:evidence_boundaries} states the paper's principal conclusions
at the level supported by their evidence. The controlled experiments remain
the main demonstration that the complete BC-assisted TD-MPC2 recipe can
complete the evaluated ForceWipe task through direct deployment. The audit
changes the interpretation of selected force-tail and disturbance claims; it
does not replace that capability result.

\begin{table*}[t]
\centering
\caption{Claim-specific evidence after the provenance and numerical audit.
Counts refer to the 225 controlled evaluations, 90 full-TD-MPC2 OOD
evaluations, 45 cross-engine evaluations, or 32 paired numerical-intervention
cells identified in each row.}
\label{tab:evidence_boundaries}
\small
\setlength{\tabcolsep}{4pt}
\begin{tabular}{@{}p{0.22\textwidth}p{0.35\textwidth}p{0.34\textwidth}@{}}
\toprule
Conclusion & Supported statement & Boundary \\
\midrule
Controlled completion & BC-assisted TD-MPC2 completed 225/225 first-pass evaluations & Exact for the frozen controlled configuration; not a physical deployment claim \\
Controlled tracking & 223/225 met the force-tracking criteria; the 12~N tier was biased low & Completion and precise regulation remain distinct \\
Weak-curvature transfer & Task and tracking passes fell to 39/90 and 37/90 & Combined block shift; individual causes are not isolated \\
Force tails & Recorded peaks and exceedances describe the baseline numerical configuration & Material peak/threshold responses occurred in 11/32 12~N OOD sentinel cells \\
MPPI contribution & Task and tracking intervals crossed zero; host planning time increased by 50.96~ms & Same-checkpoint comparison under the frozen configuration; no general peak-force claim \\
Noise robustness & No claim retained for the four unverified reference/sensor-noise labels & Declared effects were not shown to reach the learner-visible path \\
Cross-engine portability & MuJoCo achieved 30/30 compound passes on flat and inclined surfaces but 0/15 on the cylinder & Zero-shot stress test under a documented tangential calibration mismatch; not solver equivalence or hardware evidence \\
\bottomrule
\end{tabular}
\end{table*}
